\documentclass[11pt,a4paper]{report}

\usepackage[utf8]{inputenc}
\usepackage[T1]{fontenc}
\usepackage{lmodern}
\usepackage{graphicx}
\usepackage{xcolor}
\usepackage{hyperref}
\usepackage{geometry}
\usepackage{booktabs}
\usepackage{listings}
\usepackage{tcolorbox}
\usepackage{tikz}
\usetikzlibrary{shapes.geometric, arrows, positioning, shadows}

\geometry{margin=1in}

\definecolor{primary}{RGB}{0, 102, 204}
\definecolor{secondary}{RGB}{102, 102, 102}
\definecolor{codebg}{RGB}{245, 245, 245}
\definecolor{cmdcolor}{RGB}{0, 128, 0}

\hypersetup{
    colorlinks=true,
    linkcolor=primary,
    filecolor=magenta,      
    urlcolor=primary,
    pdftitle={XDP-L2-Guard Developer Guide},
    pdfpagemode=FullScreen,
}

\lstset{
    backgroundcolor=\color{codebg},
    basicstyle=\ttfamily\small,
    breaklines=true,
    frame=single,
    rulecolor=\color{secondary},
    keywordstyle=\color{primary},
    commentstyle=\color{secondary},
    stringstyle=\color{cmdcolor},
}

\title{\textbf{XDP-L2-Guard Developer Guide}}
\author{Development Team}
\date{\today}

\begin{document}

\maketitle
\tableofcontents

\chapter{Introduction}
This guide is intended for developers who wish to modify or extend XDP-L2-Guard. It covers environment setup, compilation, extension patterns, and testing procedures.

\chapter{Development Environment}

The project is optimized for Ubuntu 22.04 LTS. Ensure you have the necessary dependencies installed by running the provided setup script:

\begin{lstlisting}[language=bash]
sudo ./scripts/setup_env.sh
\end{lstlisting}

\chapter{Compilation}

The project uses a \texttt{Makefile} to compile the eBPF C code into an object file using \texttt{clang}.

\begin{lstlisting}[language=bash]
make
\end{lstlisting}

This will generate \texttt{src/data\_plane/filter.o}. The \texttt{Makefile} uses \texttt{bpftool} to generate \texttt{vmlinux.h} if needed, supporting the CO-RE (Compile Once - Run Everywhere) architecture.

\chapter{Adding New Actions}

To add a new action (e.g., \texttt{LOGGING} or \texttt{RATE\_LIMITING}), follow these steps:

\section{Update Data Plane}
Modify \texttt{src/data\_plane/filter.c}:
\begin{enumerate}
    \item Add the new action to \texttt{enum xdp\_action\_target}.
    \item Implement the logic inside the \texttt{switch (cfg->target)} block in \texttt{xdp\_drop\_logic}.
\end{enumerate}

\section{Update Control Plane}
Modify the management tools:
\begin{enumerate}
    \item \textbf{Python CLI}: Modify \texttt{src/control\_plane/guard.py} to support the new target.
    \item \textbf{Bash Wrapper}: If using the bash version, update \texttt{scripts/guard.sh}.
\end{enumerate}

\chapter{Testing \& Validation}

Several scripts are provided in the \texttt{scripts/} directory for testing different aspects of the system.

\section{Simple Ping Test}
Located in \texttt{scripts/simple\_ping/}, this set of scripts sets up a test environment using network namespaces.
\begin{itemize}
    \item \texttt{setup.sh}: Creates namespaces and virtual ethernet pairs.
    \item \texttt{ping.sh}: Runs a ping test between namespaces.
    \item \texttt{block.sh} / \texttt{unblock.sh}: Demonstrates adding/removing drop rules.
\end{itemize}

\section{Action Showcase}
Located in \texttt{scripts/action\_showcase/}, these scripts demonstrate advanced actions like \texttt{REDIRECT} and \texttt{NAT}.
\begin{itemize}
    \item \texttt{setup.sh}: Sets up a complex topology with multiple namespaces.
    \item \texttt{monitor.sh}: Runs \texttt{tcpdump} in background namespaces to verify packet arrival.
\end{itemize}

\section{Performance Testing}
Use \texttt{scripts/pktgen\_flood/pktgen\_flood.sh} to generate high volumes of UDP traffic to test the filter's performance under load.

\chapter{Debugging}

Since eBPF runs in the kernel, debugging is primarily done via \texttt{bpf\_printk}. You can view the output using:

\begin{lstlisting}[language=bash]
sudo cat /sys/kernel/debug/tracing/trace_pipe
\end{lstlisting}

Alternatively, use \texttt{bpftool} to inspect maps:

\begin{lstlisting}[language=bash]
sudo bpftool map dump name action_map
\end{lstlisting}

\end{document}
